%% ============================================================================
%%  applications/template.tex  —  SKELETON RESUME (single source template)
%% ----------------------------------------------------------------------------
%%  This is the ONLY resume template. The orchestrator copies this file into a
%%  position folder, then the writer fills the Experience / Projects / Skills
%%  sections from the context.md bullet pool. It replaces the old per-track
%%  templates (full-stack / ai-ml / dev-ops) — track is now a doctrine-routing
%%  label, not a file choice, so any track (including robotics/autonomy) fills
%%  this same skeleton.
%%
%%  WHAT IS FIXED (do not change): header (\name, \contact) and Education.
%%  WHAT GETS FILLED: Experience entries, Projects entries, Technical Skills.
%%
%%  rSectionEntry signature — 5 positional params, ALWAYS supply all 5 braces:
%%     {1 name}{2 dates}{3 tagline/role/link}{4 technologies}{5 location}
%%     - param 4 (technologies) renders after the name with a " | " separator;
%%       leave it {} for experiences, fill it for projects (the tech-stack line).
%%     - any unused param is an empty {}.
%%  rSet signature — {1 bucket label}{2 comma-separated items}.
%%
%%  PATH NOTE: this template sits at resume/applications/template.tex, beside
%%  template.cls, so the class path is just `template`. On copy into
%%  applications/{year}/{company}/{role}/, the orchestrator rewrites it to
%%  ../../../template (three levels up to applications/, where template.cls lives).
%% ============================================================================

\documentclass[
	11pt
]{../../../template}

\name{Ankur Desai}
\contact{(248) $\hspace{-0.1em}$ 657 $\hspace{-0.4em}$ -- $\hspace{-0.4em}$ 3805 \\ ardusa05@gmail.com \\ \href{https://ardusa.dev/}{www.ardusa.dev} \\ \href{https://linkedin.com/in/ardusa/}{linkedin.com/in/ardusa}}

\begin{document}

	% ---- EDUCATION (fixed; writer does not modify) --------------------------
	\begin{rSection}{E}{ducation}
		\begin{rSectionEntry}
			{Michigan State University}
			{Expected May 2028}
			{B.S. in Computer Science}
			{}
			{East Lansing, MI}

			\item \textbf{GPA:} 3.5 / 4.0
			\item \textbf{Coursework:} Data Structures \& Algorithms, OOP, Computer Architecture, Linear Algebra, Discrete Math
		\end{rSectionEntry}
	\end{rSection}

	% ---- EXPERIENCE (writer fills from context.md pool) ---------------------
	% Order by impact density for THIS role, not chronology (see resume.md I.3).
	% First bullet of each entry is the hook: lead with the problem/outcome,
	% not the technology (see resume.md I.4). 2-3 bullets per experience.
	\begin{rSection}{E}{xperience}
		\begin{rSectionEntry}
			{MindMosaic}
			{Aug 2025 -- May 2026}
			{Software Engineering Intern}
			{}
			{East Lansing, MI}

			\item Migrated hot read paths off recursive SQL self-joins onto a Neo4j graph layer so relationships stay first-class, cutting query latency 48\% (850ms to 440ms) and making multi-hop traversals 3x faster than the equivalent SQL.
			\item Built and shipped an interactive knowledge-graph platform that keeps dense, interconnected content navigable as it grows, delivering a production MVP to real users on a polyglot Java Spring Boot, Neo4j, and PostgreSQL stack.
			\item Authored an AWS deployment script and managed the environment as code with Terraform, replacing a manual 30-minute multi-service bring-up with a single-command, version-controlled provision.
		\end{rSectionEntry}

		\begin{rSectionEntry}
			{Secure and Efficient Autonomous Systems Lab}
			{May 2026 -- Present}
			{Undergraduate Research Assistant}
			{}
			{East Lansing, MI}

			\item Ran the node graph byte-identical in Gazebo and on a physical RoboMaster EP Core, driving the real robot over a custom C++ WiFi driver at a 30 Hz control loop with sub-20ms command latency.
			\item Built the verifier as a C++ ROS2 node running per-step geometric and semantic checks on each motion plan, integrating nav2's MPPI controller rather than hand-rolling one to keep the C++ work on security, not motion control.
		\end{rSectionEntry}

		\begin{rSectionEntry}
			{Claude Builder Club @ MSU}
			{Oct 2025 -- Present}
			{Co-Founder \& Vice President}
			{}
			{East Lansing, MI}

			\item Co-founded and grew an MSU engineering community from zero to 150+ members, replacing a sprawl of spreadsheets, Slack threads, and manual invites with a single club-operations platform I designed and built.
			\item Automated member onboarding so an accepted applicant is provisioned into the right GitHub teams, Slack channels, Discord roles, and Google Drive access in a single flow, cutting onboarding from 15 minutes to seconds.
			\item Tuned React Query cache tiers by data volatility across the dashboards, caching stable roles longer than fast-changing application queues to cut redundant refetches while keeping live data fresh.
		\end{rSectionEntry}
	\end{rSection}

	% ---- PROJECTS (writer fills from context.md pool) -----------------------
	% 1-2 bullets each. Tech-stack line goes in param 4 (derived from the bullets
	% selected for this project). Optional live link goes in param 3 as \href.
	\begin{rSection}{P}{rojects}
		\begin{rSectionEntry}
			{Dadei}
			{\href{https://dadei.app}{dadei.app}}
			{Ambient voice assistant that turns overheard conversation into human-approved calendar, email, and task actions}
			{React, FastAPI, Python, Redis}
			{}

			\item Built a React client that streams microphone audio up over a WebSocket with client-side Web Audio processing, receiving incremental live captions back in under 300ms as the backend transcribes each utterance.
			\item Trained a custom wake-word model to 90% recall at under one false accept per hour, gating the command pipeline so the assistant only opens the mic on demand instead of continuously streaming audio.
			\item Decoupled background workers from the live FastAPI service behind a Redis Streams event bus with consumer groups, dead-letter routing after 5 retries, and stale-claim recovery, giving at-least-once delivery that pub/sub dropped.
		\end{rSectionEntry}

		\begin{rSectionEntry}
			{fliks}
			{\href{https://github.com/fliks-gg}{github.com/fliks-gg}}
			{Multi-game platform that crowd-validates gameplay skill and mints shareable certification cards}
			{Go, PostgreSQL}
			{}

			\item Built a self-hosted video transcoding pipeline with a concurrent Go worker that pulls upload jobs, transcodes, and writes processed video back to object storage, owning the media path end to end instead of renting a managed service.
			\item Coordinated transcode jobs through a Postgres-backed queue using SELECT FOR UPDATE SKIP LOCKED instead of bolting on Redis or SQS, keeping the whole system on one datastore and one self-provisioned EC2 box.
		\end{rSectionEntry}
	\end{rSection}

	% ---- TECHNICAL SKILLS (writer picks 3-5 buckets from context.md) --------
	% ~80 char limit per line after the label. Trim tail items to fit; promote a
	% tail item only if it is a hard ATS keyword from the JD (see resume.md I.5).
	\begin{rSection}{T}{echnical Skills}
		\begin{rSet}{Languages}{Python, Go, JavaScript, TypeScript, Java, C++, SQL}
		\end{rSet}
		\begin{rSet}{Frameworks}{React, FastAPI, Spring Boot}
		\end{rSet}
		\begin{rSet}{Databases}{PostgreSQL, Neo4j}
		\end{rSet}
		\begin{rSet}{Robotics}{ROS2, Gazebo}
		\end{rSet}
	\end{rSection}

\end{document}
